\section*{Chapter \ref{ch:g10:chemistry-of-life} --- The Chemical Makeup of Living Things}

\begin{solution}{exo:g10:chemistry-of-life:1}
Oxygen (about 65\%), carbon (18.5\%), hydrogen (9.5\%) and nitrogen
(3.2\%): 96\% of body mass together.
\end{solution}

\begin{solution}{exo:g10:chemistry-of-life:2}
Mineral: water, sodium chloride, calcium phosphate, carbon dioxide (a
carbon compound, but made outside life and without a carbon--hydrogen
skeleton). Organic: glucose, the triglyceride, the \omterm{def:g10:chemistry-of-life:families}{amino acid}.
\end{solution}

\begin{solution}{exo:g10:chemistry-of-life:3}
Water lost $12.0 - 0.6 = \qty{11.4}{g}$; $11.4/12.0 = 0.95$: the leaf
is 95\% water.
\end{solution}

\begin{solution}{exo:g10:chemistry-of-life:4}
(a) Iodine solution on the crushed seed: blue-black. (b) Biuret
reagent: violet. (c) Crush the seed on paper: a translucent spot that
does not dry; Sudan III stains it red.
\end{solution}

\begin{solution}{exo:g10:chemistry-of-life:5}
\omterm{def:g10:chemistry-of-life:families}{Carbohydrates}: simple sugars such as glucose. \omterm{def:g10:chemistry-of-life:families}{Lipids}: glycerol and
fatty acids. \omterm{def:g10:chemistry-of-life:families}{Proteins}: \omterm{def:g10:chemistry-of-life:families}{amino acids} (twenty kinds). \omterm{def:g10:chemistry-of-life:families}{Nucleic acids}:
nucleotides (four kinds).
\end{solution}

\begin{solution}{exo:g10:chemistry-of-life:6}
Carbon: $18.5/0.02 \approx 900$ times more concentrated. Nitrogen:
$3.2/0.002 = 1600$ times. Living matter is not a sample of its
surroundings: it gathers a few rare elements by factors of a thousand,
which is a chemical signature of life.
\end{solution}

\begin{solution}{exo:g10:chemistry-of-life:7}
Water: $0.62 \times 70 = \qty{43.4}{kg}$. Dry matter: $70 - 43.4 =
\qty{26.6}{kg}$. \omterm{def:g10:chemistry-of-life:families}{Protein}: $0.17 \times 70 = \qty{11.9}{kg}$.
\end{solution}

\begin{solution}{exo:g10:chemistry-of-life:8}
\omterm{def:g10:chemistry-of-life:families}{Carbohydrates} $60 \times 17 = \qty{1020}{kJ}$; \omterm{def:g10:chemistry-of-life:families}{lipids} $12 \times 38 =
\qty{456}{kJ}$; \omterm{def:g10:chemistry-of-life:families}{proteins} $8 \times 17 = \qty{136}{kJ}$; total
\qty{1612}{kJ}, about \qty{1600}{kJ}. \omterm{def:g10:chemistry-of-life:families}{Lipids} supply $456/1612 \approx
28\%$ of the energy with only 12\% of the mass.
\end{solution}

\begin{solution}{exo:g10:chemistry-of-life:9}
Bread: blue-black --- flour is mostly starch. Cheese: no colour change
beyond the iodine's own tint --- \omterm{def:g10:chemistry-of-life:families}{proteins} and \omterm{def:g10:chemistry-of-life:families}{lipids}, no starch. Sugar:
no change --- sucrose is a simple sugar, not starch; iodine detects the
long starch chain only.
\end{solution}

\begin{solution}{exo:g10:chemistry-of-life:10}
The dry matter is mostly organic: carbon--hydrogen molecules that
combine with oxygen, releasing energy, carbon dioxide and water. The
ash is the \omterm{def:g10:chemistry-of-life:mineral-organic}{mineral salts}, already fully combined with oxygen; nothing
in them can burn.
\end{solution}

\begin{solution}{exo:g10:chemistry-of-life:11}
Fehling's reagent reacts with a free chemical group carried by single
glucose molecules; in starch the glucose units are joined to one
another and those groups are locked inside the chain, so the reagent
finds nothing. Iodine, conversely, lodges in the long coiled starch
chain, which single glucose molecules do not form. Each test detects a
structure, not an element.
\end{solution}

\begin{solution}{exo:g10:chemistry-of-life:12}
The grain has absorbed water massively. At 13\% the reserves are solid
and the enzymes cannot work: the grain's chemistry is switched off and
it survives dormant. Water dissolves the stores, lets the enzymes act
and carries the products to the growing embryo: chemistry restarts
only once water is abundant.
\end{solution}

\begin{solution}{exo:g10:chemistry-of-life:13}
Salt: $9 \times 0.5 = \qty{4.5}{g}$; glucose: $1 \times 0.5 =
\qty{0.5}{g}$. Pure water would dilute the plasma; water would then
enter the red blood cells, which would swell and burst.
\end{solution}

\begin{solution}{exo:g10:chemistry-of-life:14}
Glycogen: $500 \times 17 = \qty{8500}{kJ}$. \omterm{def:g10:chemistry-of-life:families}{Lipid}: $10000 \times 38 =
\qty{3.8e5}{kJ}$. To hold \qty{3.8e5}{kJ} as glycogen: $380000/17
\approx \qty{22.4}{kg}$ of glycogen, plus three times that mass of
water, about \qty{90}{kg} in all --- more than the person. Fat stores
roughly nine times more energy per kilogram carried, which is why
animals that must move store their reserves as \omterm{def:g10:chemistry-of-life:families}{lipid}.
\end{solution}

\begin{solution}{exo:g10:chemistry-of-life:15}
The atoms are indeed ordinary, but (1) the proportions are not: living
matter concentrates carbon a thousandfold and nitrogen far more,
against an oxygen--silicon crust; (2) the molecules are not: only
living things assemble carbon into giant chains --- \omterm{def:g10:chemistry-of-life:families}{proteins}, starch,
DNA --- whose size has no mineral counterpart; (3) those molecules
carry information in the order of their units, and no crystal does.
Life is chemically special in its organisation, not in its atoms.
\end{solution}

\begin{solution}{pb:g10:chemistry-of-life:1}
\textbf{1.} Water $0.13 \times 45 = \qty{5.9}{mg}$; dry mass
\qty{39.2}{mg}.

\textbf{2.} Starch $0.70 \times 39.2 \approx \qty{27.4}{mg}$; \omterm{def:g10:chemistry-of-life:families}{protein}
$0.12 \times 39.2 \approx \qty{4.7}{mg}$; \omterm{def:g10:chemistry-of-life:families}{lipid} $0.02 \times 39.2
\approx \qty{0.8}{mg}$.

\textbf{3.} Iodine (blue-black) for starch; biuret (violet) for
\omterm{def:g10:chemistry-of-life:families}{protein}; a translucent spot on paper or Sudan III for \omterm{def:g10:chemistry-of-life:families}{lipid}.

\textbf{4.} At 13\% the stores are solid and the enzymes inactive;
water must dissolve the reserves and let the reactions run before the
embryo can grow. Germination starts by absorbing water, and the
seedling's 88\% is the sign that its chemistry is on.

\textbf{5.} Water $0.13 \times 8 = \qty{1.04}{t}$; dry matter
\qty{6.96}{t}, about \qty{7}{t}.

\textbf{6.} Mass $500 + 320 + 10 - 130 = \qty{700}{g}$. Water in the
loaf: from the flour $0.13 \times 500 = \qty{65}{g}$, plus \qty{320}{g},
minus \qty{130}{g}: \qty{255}{g}, i.e. $255/700 \approx 36\%$.

\textbf{7.} Flour dry matter $0.87 \times 500 = \qty{435}{g}$; starch
$0.70 \times 435 \approx \qty{305}{g}$; \omterm{def:g10:chemistry-of-life:families}{protein} $0.12 \times 435 \approx
\qty{52}{g}$ (\omterm{def:g10:chemistry-of-life:families}{lipid} about \qty{9}{g}).

\textbf{8.} $305 \times 17 + 52 \times 17 + 9 \times 38 \approx 5185 +
884 + 342 \approx \qty{6400}{kJ}$.

\textbf{9.} $120/700 \approx 0.17$ of the loaf: about \qty{1100}{kJ},
i.e. 11\% of \qty{10000}{kJ}.

\textbf{10.} Salt is mineral: it is already fully oxidised and has no
carbon--hydrogen bonds to burn. In the body its ions set the salinity
of blood and are needed for nerve and muscle signals.

\textbf{11.} $0.60 \times 60 = \qty{36}{kg}$ of water.

\textbf{12.} \omterm{def:g10:chemistry-of-life:families}{Protein} $0.17 \times 60 = \qty{10.2}{kg}$; \omterm{def:g10:chemistry-of-life:families}{lipid} $0.16
\times 60 = \qty{9.6}{kg}$; minerals \qty{3.6}{kg}; \omterm{def:g10:chemistry-of-life:families}{carbohydrates}
\qty{0.6}{kg}.

\textbf{13.} $1.2/36 \approx 3.3\%$ of her water. Blood volume and
temperature control both depend on it; a loss of 10\% is dangerous, so
drinking must restore it within hours.

\textbf{14.} Plasma $0.55 \times 4.5 \approx \qty{2.5}{L}$; water $0.9
\times 2.5 \approx \qty{2.2}{kg}$, about 6\% of body water. The rest is
inside the cells (about two thirds of all body water) and in the fluid
between cells.

\textbf{15.} An infant exchanges a far larger fraction of its water
each day (drinking, urine, skin) than an adult, so a few hours of loss
without intake removes a dangerous fraction of its total.

\textbf{16.} \omterm{def:g10:chemistry-of-life:families}{Carbohydrates} $0.40 \times 0.6 = \qty{0.24}{kg}$; \omterm{def:g10:chemistry-of-life:families}{lipids}
$0.77 \times 9.6 \approx \qty{7.4}{kg}$; \omterm{def:g10:chemistry-of-life:families}{proteins} $0.53 \times 10.2
\approx \qty{5.4}{kg}$. Total about \qty{13}{kg} of carbon.

\textbf{17.} $13/60 \approx 22\%$, against 18.5\%. The agreement is
good given rounded family fractions and the variation of \omterm{def:g10:chemistry-of-life:families}{lipid} content
between individuals; a leaner body gives a lower figure.

\textbf{18.} $0.16 \times 10.2 \approx \qty{1.6}{kg}$ of nitrogen. It
came from the nitrogen of the air, captured by soil bacteria, taken up
by plants and eaten along the food chain.

\textbf{19.} $0.89 \times 36 \approx \qty{32}{kg}$ of oxygen, i.e.
$32/60 \approx 53\%$ of body mass. Water alone accounts for most of the
65\%; the remainder is the oxygen inside the organic molecules.

\textbf{20.} A \qty{60}{kg} person carries about \qty{13}{kg} of
carbon. It left the air as carbon dioxide through photosynthesis in
plants, and entered the body through feeding and digestion along the
food chain.
\end{solution}
